\chapter*{Použité technologie}
\addcontentsline{toc}{chapter}{Použité technologie}

Vzhledem k rozsahu projektu bylo využito široké spektrum technologií od nízkoúrovňového programování
až po mechanický návrh. Výběr nástrojů reflektoval technické požadavky, dostupnost dokumentace a
také preferenci otevřeného softwaru (\textit{FOSS}). Klíčové technologie jsou shrnuty v tabulce
\ref{tab:technologies}.

\begin{table}[htbp]
    \centering
    \small
    % \renewcommand{\arraystretch}{0.75}
    \begin{tabularx}{\textwidth}{l l >{\raggedright\arraybackslash}X}
        \toprule
        \textbf{Kategorie} & \textbf{Technologie} & \textbf{Využití v projektu} \\
        \midrule

        \textbf{Vývoj a sestavení}
        & Jazyk C, GCC
        & Hlavní implementační jazyk a kompilátor pro architekturu ARM (RP2350). \\
        & CMake
        & Automatizace sestavení, správa závislostí a multiplatformní podpora. \\
        & Git
        & Správa verzí, historie změn a synchronizace zdrojového kódu. \\
        \addlinespace

        \textbf{Emulace a BIOS}
        & cc65
        & Toolchain pro 6502; sestavení BIOSu a testovacích programů. \\
        & neslib
        & Vývoj UI BIOSu (práce s registry PPU, správa palet a textu). \\
        & Unity
        & Unit testing framework pro validaci implementace instrukcí CPU 6502. \\
        \addlinespace

        \textbf{Embedded systém}
        & Pico SDK
        & Oficiální abstrakční vrstva (HAL) pro nízkoúrovňové řízení RP2350. \\
        & PicoDVI
        & Softwarové generování DVI/HDMI signálu pomocí PIO a DMA. \\
        & FatFS
        & Implementace souborového systému FAT32 pro načítání dat z SD karty. \\
        \addlinespace

        \textbf{Hardware a kryty}
        & KiCad
        & Návrh schémat a desky plošných spojů (PCB) emulátoru. \\
        & Fusion 360
        & Konstrukce a 3D modelování mechanické schránky zařízení. \\
        & PrusaSlicer
        & Příprava dat (G-code) pro 3D tisk součástí krytu. \\
        \addlinespace

        \textbf{Desktop backend}
        & SDL3
        & Abstrakce grafiky, zvuku a vstupů pro ladění emulačního jádra na PC. \\
        \addlinespace

        \textbf{Dokumentace}
        & {\fontfamily{cmr}\selectfont \hologo{LaTeX}}
        & Sazba textu této práce. \\
        & {\fontfamily{cmr}\selectfont \hologo{LuaLaTeX}}
        & Kompilátor zdrojových souborů textové práce v sázecím systému
        {\fontfamily{cmr}\selectfont \hologo{LaTeX}}. \\
        & PlantUML
        & Kreslení časových diagramů. \\
        & draw.io
        & Kreslení diagramů a ilustrací. \\

        \bottomrule
    \end{tabularx}
    \caption{Přehled použitých technologií, knihoven a nástrojů.}
    \label{tab:technologies}
\end{table}
